TODO GERAL:

- Corrigir o fortmato de todas as citações
- Corrigir as Legendas de todas as figuras



% II. Referencial Bibliografico
\section{Referencial Bibliografico}

\subsection{Modelos de Linguagem e a Arquitetura de Harness}
- Explciar os thinking e oq os principais trabalhos quie trazem o termo Harness abordam

\subsection{Degradação de Contexto}
- Mostrar o contraponto do thinking e hardness

\subsection{O Ralph Wiggum Loop}
-


TODO:
    - Analisar tudo oque tem aq
    - Preciso garantir que tenha pelo menos a explicação do Ralph Loop. Nessa explicação precisa ficar muito claro o porque dele funcionar somente em casos de engenharia de software


ORIGINAL RALPH WHIGGLE LOOP:

- Explicar todo o passo a passo, mas dar foco no "iMPLEMENTA A TEREFA"
    "Fresh context each loop", "No context compaction", "Backpressure via tests"
    - Explicar que essa fase tb tem loop, aq são definidos testes unitarios e ate que esses testes passem , a historia nao é dada como concluida
    - Deixar claro oq é stateless e oq nao


% ==================================================================== %
% III. Metodologia
\section{Metodologia}
\titleformat{\subsection}
  {\color{grayFEI}\normalfont\fontfamily{phv}\selectfont\itshape\fontsize{10}{12}\selectfont}
  {\arabic{subsection}.}        % Numeração arábica
  {0.2em}
  {}
\renewcommand{\thesubsection}{\arabic{subsection}}  % Define arábica globalmente
\titlespacing*{\subsection}{0pt}{6pt}{3pt}

TODO:
    - Colocar a visão geral de como queremos adaptar o ralph loop para que ele funcione para responder perguntas gerais:
        - O documento de reuquisitos vira um input generica, como uma pergunta.
        - As historias de usuario que originalmente são criadas a partir dos requisitos, agora sã otarefas que são escritas a partir da pergunta do usuario. Seguindo a mesma logica anterior, cada tarefa tem seus proprios criterios de aceite
        - Aqui preciso deixar claro que se for dar exemplos vou usar JSON, mas isso é so por conviniencia, e poderia ser qualquer outra maneira que possa armazenar texto e pegar partes dele mais tarde


% 1. Definição dos critérios de aceite para a resposta final
\subsection{Definição dos critérios de aceite para a resposta final}
-


% 2. Quebra de tarefas independentes
\subsection{Quebra de tarefas independentes}



% 3. Resolução das tarefas
\subsection{Resolução das tarefas}



% 4. Compara do resultado compilado das tarefas
\subsection{Compara do resultado compilado das tarefas}


